\documentclass[11pt,a4paper]{article}

\usepackage[utf8]{inputenc}
\usepackage[T1]{fontenc}
\usepackage[french]{babel}

\usepackage{amsmath,amssymb,amsthm}

\usepackage{geometry}
\geometry{margin=2.5cm}

\usepackage{array}
\usepackage{longtable}
\usepackage{booktabs}

\usepackage{listings}
\usepackage{xcolor}

\definecolor{codebg}{rgb}{0.96,0.96,0.96}
\definecolor{codeframe}{rgb}{0.8,0.8,0.8}

\lstset{
  backgroundcolor=\color{codebg},
  frame=single,
  rulecolor=\color{codeframe},
  basicstyle=\ttfamily\small,
  breaklines=true,
  columns=fullflexible,
  keepspaces=true,
  showstringspaces=false
}

\newcommand{\sem}[4]{%
  \langle #1,\;#2\rangle \rightarrow \langle #3,\;#4\rangle}
\newcommand{\ext}[3]{#1\!\left[#2 \mapsto #3\right]}

\newcommand{\irule}[2]{%
  \dfrac{\displaystyle #1}{\displaystyle #2}
}

\newcommand{\dom}[1]{\mathrm{dom}(#1)}

\newcommand{\truthy}{\mathrm{truthy}}

\newcommand{\V}[1]{\mathtt{#1}}

\newenvironment{bloc}[1]
{
  \par\medskip
  \noindent\textbf{#1}
  \par
  \smallskip
  \begin{quote}
}
{
  \end{quote}
  \medskip
}

\title{
\textbf{Résumé Projet IN213}\\[0.5em]
\large CaioSHell --- Interpréteur de shell en OCaml
}

\author{
Caio ESCORCIO LIMA DOURADO\\
ENSTA Paris
}

\date{25 Mai 2026}


\begin{document}


\maketitle
\tableofcontents
\clearpage

% ============================================================
\section{Vue d'ensemble}
% ============================================================

\textbf{Objectif :}
concevoir un interpréteur d'un langage shell personnalisé
(\texttt{.csh}), exécuté directement via les appels système Unix,
sans compilation préalable.

\subsection{Architecture}

\begin{longtable}{|p{4cm}|p{9cm}|}
\hline
\textbf{Fichier} & \textbf{Rôle} \\
\hline
\texttt{cshlex.mll} & Analyse lexicale (tokens) \\
\hline
\texttt{cshparse.mly} & Analyse syntaxique (AST) \\
\hline
\texttt{cshast.ml} & Définition des nœuds AST \\
\hline
\texttt{cshinterp.ml} & Interpréteur / évaluateur \\
\hline
\texttt{cshloop.ml} & Boucle REPL interactive \\
\hline
\texttt{test.csh} & Exemples de la syntaxe du code \\
\hline
\end{longtable}

\subsection{Pipeline d'exécution}

\begin{lstlisting}
                                source .csh
                                    V
                                cshlex.mll
                                    V
                                  tokens
                                    V
                                cshparse.mly
                                    V
                                   AST
                                    V
                                cshinterp.ml
                                    V
                                fork / execvp
                                    V
                                resultat Unix
\end{lstlisting}

\begin{bloc}{Exemple}
\begin{lstlisting}
ls -la /tmp;;
\end{lstlisting}

\begin{itemize}
  \item Lexer :
\begin{lstlisting}
IDENT PATH PATH SEMISEMI
\end{lstlisting}

  \item Parser :
\begin{lstlisting}
ECmd[
  EIdent "ls";
  EIdent "-la";
  EIdent "/tmp"
]
\end{lstlisting}

  \item Interpréteur :
\begin{lstlisting}
fork + execvp "ls"
[|"ls"; "-la"; "/tmp"|]
\end{lstlisting}
\end{itemize}
\end{bloc}

% ------------------------------------------------------------
\subsection{Inspiration générale}
% ------------------------------------------------------------

L'implémentation de \texttt{cshinterp.ml} s'appuie sur deux sources
d'inspiration principales.

\medskip
\textbf{\texttt{foo.ml} --- modèle pour les appels système Unix.}
Le fichier \texttt{foo.ml} fourni dans le cours définit deux fonctions,
\texttt{exec\_command} et \texttt{pipe\_commands}, qui illustrent
l'utilisation de \texttt{Unix.fork}, \texttt{Unix.execvp},
\texttt{Unix.pipe} et \texttt{Unix.waitpid}.
Les fonctions \texttt{eval\_cmd} et \texttt{eval\_pipe} reprennent
directement cette structure : séparation fils/père, nommage des PIDs
(\texttt{child\_pid}, \texttt{dead\_pid}), flag \texttt{Unix.WUNTRACED}
et gestion des codes retour.

\medskip
\textbf{Module \texttt{List} d'OCaml --- fonctions réimplémentées.}
L'interpréteur \emph{n'importe pas} le module \texttt{List}.
Les fonctions nécessaires (\texttt{map}, \texttt{fold\_left},
\texttt{fold\_left2}, \texttt{iter}, \texttt{assoc}, \texttt{assoc\_opt})
sont réécrites à la main, ce qui rend leur rôle explicite dans le
contexte de l'interpréteur. De même, \texttt{to\_array} remplace
\texttt{Array.of\_list} sans faire appel à la bibliothèque standard.

\subsection{Typage}

Le langage possède un \textbf{typage dynamique} :
aucune vérification statique n'est effectuée avant l'exécution.

Les erreurs de types apparaissent uniquement au moment
où les valeurs sont réellement utilisées.


% ============================================================
\section{Valeurs et environnement}
% ============================================================

\subsection{Le type \texttt{value}}

\begin{lstlisting}[language=ml]
type value =
  | VInt    of int
  | VBool   of bool
  | VString of string
  | VUnit
  | VClosure of (string list * Cshast.expr * env)
\end{lstlisting}

\begin{center}
\begin{tabular}{|p{4cm}|p{8cm}|}
\hline
\textbf{Constructeur} & \textbf{Exemple} \\
\hline
\texttt{VInt 42} & \texttt{42}, \texttt{@x + 1} \\
\hline
\texttt{VBool true} & booléens, codes retour Unix \\
\hline
\texttt{VString "abc"} & chaînes, chemins, commandes \\
\hline
\texttt{VUnit} & absence de valeur significative \\
\hline
\texttt{VClosure} & fonction utilisateur \\
\hline
\end{tabular}
\end{center}

\begin{bloc}{Closures}
Une \texttt{VClosure} capture :

\begin{enumerate}
  \item les paramètres,
  \item le corps AST,
  \item l'environnement $\sigma_f$ au moment de la définition.
\end{enumerate}

Cela implémente la \textbf{portée lexicale}.
\end{bloc}

\subsection{L'environnement $\sigma$}

\begin{lstlisting}[language=ml]
type env = (string * value) list
\end{lstlisting}

L'environnement est une liste d'association :

\[
\sigma =
[
("x", \V{VInt}\ 10);
("name", \V{VString}\ "Alice")
]
\]

\subsection{Extension de l'environnement}

\begin{lstlisting}[language=ml]
let bind x v env = (x,v) :: env
\end{lstlisting}

Notation sémantique :

\[
\ext{\sigma}{x}{v}
\]

\begin{bloc}{Masquage (shadowing)}
\[
\sigma = [(x,5)]
\qquad
\ext{\sigma}{x}{10} = [(x,10);(x,5)]
\]

La recherche retourne toujours la liaison la plus récente.
Un ancien lien n'est jamais supprimé --- il est simplement masqué.
\end{bloc}

% ============================================================
\section{Sémantique opérationnelle}
% ============================================================

\subsection{Jugement d'évaluation}

\[
\sem{e}{\sigma}{v}{\sigma'}
\]

Lecture :

\begin{quote}
« L'expression $e$, évaluée dans l'environnement $\sigma$,
produit la valeur $v$ et le nouvel environnement $\sigma'$. »
\end{quote}

Correspondance directe avec OCaml :

\begin{lstlisting}[language=ml]
val eval :
  Cshast.expr -> env -> value * env
\end{lstlisting}

\subsection{Règles d'inférence}

Forme générale :

\[
\irule{\text{prémisses}}{\text{conclusion}}
\]

Une règle sans prémisse est un axiome.

% ============================================================
\section{Règles d'évaluation}
% ============================================================

\subsection{Littéraux}

\[
\sem{\V{EInt}\;n}{\sigma}{\V{VInt}\;n}{\sigma}
\qquad
\sem{\V{EBool}\;b}{\sigma}{\V{VBool}\;b}{\sigma}
\]

\[
\sem{\V{EString}\;s}{\sigma}{\V{VString}\;s}{\sigma}
\qquad
\sem{\V{EPass}}{\sigma}{\V{VUnit}}{\sigma}
\]

Les littéraux ne modifient jamais l'environnement.
\texttt{EPass} est un no-op syntaxique utilisé dans les branches vides.

\subsection{\texttt{EIdent} et \texttt{EVar}}

\textbf{Identifiant libre (\texttt{EIdent}) :}

\[
\irule{
s \in \dom{\sigma}
}{
\sem{\V{EIdent}\;s}{\sigma}{\sigma(s)}{\sigma}
}
\qquad
\irule{
s \notin \dom{\sigma}
}{
\sem{\V{EIdent}\;s}{\sigma}{\V{VString}\;s}{\sigma}
}
\]

\begin{bloc}{Interprétation}
Si l'identifiant n'existe pas dans $\sigma$,
il devient automatiquement une chaîne de caractères.
Cela permet d'écrire \texttt{ls /tmp} sans guillemets :
\texttt{"ls"} et \texttt{"/tmp"} sont absents de $\sigma$
et deviennent des \texttt{VString} passés à \texttt{execvp}.
\end{bloc}

\textbf{Variable stricte (\texttt{EVar}, syntaxe \texttt{@x}) :}

\[
\irule{
s \in \dom{\sigma}
}{
\sem{\V{EVar}\;s}{\sigma}{\sigma(s)}{\sigma}
}
\qquad
s \notin \dom{\sigma}
\Rightarrow
\texttt{Unbound variable}
\]

L'absence d'une variable \texttt{@x} est une erreur de programmation,
contrairement à un identifiant libre.

\subsection{\texttt{ESeq}}

\[
\irule{
\sem{e_1}{\sigma}{v_1}{\sigma_1}
\quad
\cdots
\quad
\sem{e_n}{\sigma_{n-1}}{v_n}{\sigma_n}
}{
\sem{\V{ESeq}[e_1;\dots;e_n]}{\sigma}{v_n}{\sigma_n}
}
\]

L'environnement est \textbf{propagé de gauche à droite} :
une définition en position $i$ est visible aux positions $i+1 \ldots n$.
Implémenté par \texttt{fold\_left} :

\begin{lstlisting}[language=ml]
let step (_, acc_env) e = eval e acc_env in
fold_left step (VUnit, env) es
\end{lstlisting}

\subsection{\texttt{EDefVar} et \texttt{EAssign}}

\[
\irule{
\sem{e}{\sigma}{v}{\sigma'}
}{
\sem{\V{EDefVar}(x,e)}
{\sigma}
{\V{VUnit}}
{\ext{\sigma'}{x}{v}}
}
\]

\begin{bloc}{Observation}
\texttt{EDefVar} (\texttt{define var x = e}) et \texttt{EAssign} (\texttt{x = e})
utilisent exactement la même règle sémantique :
la différence est \textbf{uniquement syntaxique}.
\end{bloc}

\subsection{\texttt{EDefFun}}

\[
\sem{
\V{EDefFun}(f,params,body)
}{
\sigma
}{
\V{VUnit}
}{
\ext{\sigma}{f}{
\V{VClosure}(params,body,\sigma)
}
}
\]

\begin{bloc}{Portée lexicale}
La fermeture capture l'environnement $\sigma$
\textbf{au moment de la définition}, pas au moment de l'appel.
Si une variable est redéfinie après \texttt{define function},
la fonction utilise toujours la valeur capturée.
\end{bloc}

\subsection{\texttt{EApp}}

\[
\irule{
\sigma(f)=\V{VClosure}(ps,body,\sigma_f)
\quad
\forall i:\sem{a_i}{\sigma}{v_i}{\_}
\quad
\sem{body}{\sigma_c}{v}{\_}
}{
\sem{\V{EApp}(f,args)}{\sigma}{v}{\sigma}
}
\]

\begin{lstlisting}[language=ml]
| Cshast.EApp (f, args) ->
    (match lookup f env with
     | VClosure (params, body, closure_env) ->
         let eval_arg a = fst (eval a env) in
         let bind_param new_env param value = bind param value new_env in
         let values   = map eval_arg args in
         let call_env =
           try fold_left2 bind_param closure_env params values
           with Invalid_argument _ ->
             failwith (f ^ ": wrong number of arguments")
         in
         let (v, _) = eval body call_env in
         (v, env)    (* l'env de l'appelant est inchange *)
     | _ -> failwith (f ^ " is not a function"))
\end{lstlisting}

\begin{bloc}{Points clés}
\begin{itemize}
  \item Les arguments sont évalués dans l'env de \textbf{l'appelant}.
  \item Le corps s'exécute dans $\sigma_f$ (env capturé) étendu des paramètres.
  \item Les modifications du corps sont \textbf{locales} : $\sigma$ de l'appelant est retourné inchangé.
  \item \texttt{fold\_left2} associe chaque paramètre à sa valeur ; si les listes diffèrent en taille, \texttt{Invalid\_argument} est levé.
\end{itemize}
\end{bloc}

\subsection{\texttt{EBinop}}

\[
\irule{
\sem{e_1}{\sigma}{v_1}{\sigma_1}
\quad
\sem{e_2}{\sigma_1}{v_2}{\sigma_2}
}{
\sem{
\V{EBinop}(op,e_1,e_2)
}{
\sigma
}{
op(v_1,v_2)
}{
\sigma_2
}
}
\]

\begin{center}
\begin{tabular}{|l|l|l|}
\hline
Opérateur & Types & Résultat \\
\hline
\texttt{+} & int $\times$ int & somme entière \\
\texttt{+} & str $\times$ str/int & concaténation \\
\texttt{-}, \texttt{*} & int $\times$ int & entier \\
\texttt{\#} & int $\times$ int & division entière \\
\texttt{=} & tout type & booléen structural \\
\texttt{> < >= <=} & int ou str & booléen \\
\hline
\end{tabular}
\end{center}

\begin{bloc}{Pourquoi \texttt{\#} pour la division ?}
Le caractère \texttt{/} est réservé par le lexeur aux chemins de
fichiers (token \texttt{PATH}). La division entière utilise donc
\texttt{\#} : \texttt{10 \# 3} $\to$ \texttt{VInt 3}.
\end{bloc}

\subsection{\texttt{EIf} et \texttt{EAnd}}

\textbf{\texttt{EIf} :}
\[
\irule{
\sem{cond}{\sigma}{v}{\sigma'} \;\; \truthy(v) \;\;
\sem{t}{\sigma'}{v_t}{\sigma_t}
}{
\sem{\V{EIf}(cond,t,f)}{\sigma}{v_t}{\sigma_t}
}
\quad
\irule{
\sem{cond}{\sigma}{v}{\sigma'} \;\; \neg\truthy(v) \;\;
\sem{f}{\sigma'}{v_f}{\sigma_f}
}{
\sem{\V{EIf}(cond,t,f)}{\sigma}{v_f}{\sigma_f}
}
\]

\textbf{\texttt{EAnd} (court-circuit) :}
\[
\irule{
\sem{e_1}{\sigma}{v_1}{\sigma_1} \;\; \truthy(v_1) \;\;
\sem{e_2}{\sigma_1}{v_2}{\sigma_2}
}{
\sem{\V{EAnd}(e_1,e_2)}{\sigma}{v_2}{\sigma_2}
}
\quad
\irule{
\sem{e_1}{\sigma}{v_1}{\sigma_1} \;\; \neg\truthy(v_1)
}{
\sem{\V{EAnd}(e_1,e_2)}{\sigma}{\V{VBool\ false}}{\sigma_1}
}
\]

Fonction \texttt{is\_truth} :
$\V{VBool}\;b \to b$,\;
$\V{VInt}\;n \to n\neq 0$,\;
$\V{VString}\;s \to s\neq\texttt{""}$,\;
$\V{VUnit} \to \texttt{false}$.

\subsection{\texttt{EWhile}}

\[
\irule{
\sem{cond}{\sigma}{v}{\sigma'} \quad \neg\truthy(v)
}{
\sem{\V{EWhile}(cond,body)}{\sigma}{\V{VUnit}}{\sigma'}
}
\;\textsc{[Stop]}
\qquad
\irule{
\sem{cond}{\sigma}{v}{\sigma'} \quad \truthy(v)
\quad \sem{body}{\sigma'}{\_}{\sigma''}
\quad \sem{\V{EWhile}(\ldots)}{\sigma''}{u}{\sigma'''}
}{
\sem{\V{EWhile}(cond,body)}{\sigma}{u}{\sigma'''}
}
\;\textsc{[Iter]}
\]

Une référence mutable \texttt{env\_ref} transporte $\sigma$ entre les
itérations, de sorte qu'une affectation dans le corps (ex.\ \texttt{i = @i + 1})
soit visible lors de la prochaine évaluation de la condition.

\subsection{\texttt{EFor}}

\[
\irule{
\sem{e_i}{\sigma}{v_i}{\_}
\qquad
\sem{body}{\ext{\sigma}{x}{v_i}}{\_}{\sigma_i}
}{
\text{itération sur } e_i : \sigma \leftarrow \sigma_i
}
\]

À chaque itération, la variable de boucle $x$ est liée à la valeur
courante dans une extension de \texttt{env\_ref}.
Le corps est réévalué pour chaque élément de la liste.

% ============================================================
\section{Fonctions utilitaires}
% ============================================================

\subsection{Réimplémentation sans le module \texttt{List}}

Par choix pédagogique, \texttt{cshinterp.ml} n'importe aucun module
de la bibliothèque standard OCaml (\texttt{List}, \texttt{Array},
\textit{etc.}). Les fonctions suivantes sont réécrites à la main :

\begin{center}
\begin{longtable}{|l|l|p{4.5cm}|l|}
\hline
\textbf{Fonction} & \textbf{Équivalent stdlib} & \textbf{Rôle} & \textbf{Utilisée dans} \\
\hline
\texttt{assoc} & \texttt{List.assoc} & Recherche par clé, lève \texttt{Not\_found} & \texttt{lookup} \\
\hline
\texttt{assoc\_opt} & \texttt{List.assoc\_opt} & Recherche par clé, retourne \texttt{option} & \texttt{EIdent} \\
\hline
\texttt{map} & \texttt{List.map} & Applique $f$ à chaque élément & \texttt{EApp}, \texttt{eval\_cmd} \\
\hline
\texttt{fold\_left} & \texttt{List.fold\_left} & Réduction séquentielle gauche$\to$droite & \texttt{ESeq} \\
\hline
\texttt{fold\_left2} & \texttt{List.fold\_left2} & Réduction parallèle sur deux listes & \texttt{EApp} \\
\hline
\texttt{iter} & \texttt{List.iter} & Parcours avec effets de bord & \texttt{EFor} \\
\hline
\texttt{to\_array} & \texttt{Array.of\_list} & Conversion liste $\to$ tableau & \texttt{eval\_cmd} \\
\hline
\end{longtable}
\end{center}

\subsection{\texttt{assoc} et \texttt{assoc\_opt} --- recherche dans l'environnement}

La recherche est linéaire et retourne la \textbf{première} (donc la plus récente)
liaison trouvée. C'est ce qui implémente le masquage (\textit{shadowing}).

\begin{lstlisting}[language=ml]
let rec assoc key = function
  | []                        -> raise Not_found
  | (k, v) :: _ when k = key -> v
  | _ :: rest                 -> assoc key rest
\end{lstlisting}

\texttt{assoc\_opt} retourne \texttt{None} au lieu de lever une exception.
Elle est utilisée par \texttt{EIdent} pour le repli sur \texttt{VString} :
si un nom n'est pas lié dans $\sigma$, il est traité comme une chaîne littérale
(utile pour les noms de commandes comme \texttt{ls}, \texttt{grep}, etc.).

\subsection{\texttt{fold\_left} --- propagation de l'environnement}

\begin{lstlisting}[language=ml]
let rec fold_left f acc = function
  | []      -> acc
  | x :: xs -> fold_left f (f acc x) xs
\end{lstlisting}

Dans \texttt{ESeq}, l'accumulateur est la paire \texttt{(valeur, env)}.
À chaque étape, la valeur précédente est ignorée (\texttt{\_}) et le
nouvel environnement est transmis à l'expression suivante :

\begin{lstlisting}[language=ml]
let step (_, acc_env) e = eval e acc_env in
fold_left step (VUnit, env) es
\end{lstlisting}

\subsection{\texttt{fold\_left2} --- liaison des paramètres}

\begin{lstlisting}[language=ml]
let rec fold_left2 f acc l1 l2 =
  match l1, l2 with
  | [],      []       -> acc
  | x :: xs, y :: ys -> fold_left2 f (f acc x y) xs ys
  | _                 -> raise (Invalid_argument "fold_left2")
\end{lstlisting}

Dans \texttt{EApp}, elle lie chaque paramètre à sa valeur dans
\texttt{closure\_env} :

\begin{lstlisting}[language=ml]
let bind_param new_env param value = bind param value new_env in
fold_left2 bind_param closure_env params values
\end{lstlisting}

Si les deux listes ne sont pas de même longueur,
\texttt{Invalid\_argument} est capturée et convertie en
\texttt{"wrong number of arguments"}.

\subsection{\texttt{to\_array} --- conversion pour \texttt{execvp}}

\texttt{Unix.execvp} attend un \texttt{string array}, mais les arguments
sont construits comme une liste. \texttt{to\_array} effectue la conversion
sans \texttt{Array.of\_list} :

\begin{lstlisting}[language=ml]
let to_array lst = match lst with
  | [] -> [||]
  | hd :: _ ->
      let rec len acc = function
        | []     -> acc
        | _ :: t -> len (acc+1) t in
      let arr = Array.make (len 0 lst) hd in
      let rec fill i = function
        | []      -> ()
        | x :: xs -> arr.(i) <- x ; fill (i+1) xs
      in fill 0 lst ; arr
\end{lstlisting}

% ============================================================
\section{Commandes Unix}
% ============================================================

\subsection{Vue générale de \texttt{eval\_cmd}}

Inspirée de \texttt{exec\_command} dans \texttt{foo.ml},
\texttt{eval\_cmd} suit trois étapes :

\begin{enumerate}
  \item \textbf{Évaluation des arguments} : chaque expression est évaluée
        puis convertie en chaîne via \texttt{value\_to\_string}.
  \item \textbf{Filtrage des commandes internes} (\textit{builtins}) :
        \texttt{cd} et \texttt{exit} sont traitées directement dans le processus shell.
  \item \textbf{Exécution externe} : toute autre commande est lancée via \texttt{fork + execvp}.
\end{enumerate}

\begin{lstlisting}[language=ml]
let arg_to_str a = value_to_string (fst (eval a env)) in
let strs = map arg_to_str args in
\end{lstlisting}

\subsection{Commandes internes (\textit{builtins})}

\begin{longtable}{|p{3cm}|p{9cm}|}
\hline
\textbf{Commande} & \textbf{Raison d'être interne} \\
\hline
\texttt{cd} & \texttt{Unix.chdir} dans un processus fils ne change
              que le répertoire de ce fils, pas celui du shell courant.
              La commande doit donc s'exécuter dans le processus shell lui-même. \\
\hline
\texttt{exit} & Doit terminer l'interpréteur lui-même.
                Un \texttt{exit} dans un fils ne terminerait que le fils. \\
\hline
\end{longtable}

\subsection{Exécution externe : \texttt{fork + execvp}}

Structure directement inspirée de \texttt{exec\_command} dans \texttt{foo.ml} :

\begin{lstlisting}[language=ml]
match Unix.fork () with
| 0 ->
    (* Fils : remplace son image par la commande a executer.
       Le premier argument de argv[] est le nom de la commande. *)
    Unix.execvp cmd (to_array (cmd :: rest))
| child_pid ->
    (* Pere : attend explicitement la fin de ce fils. *)
    match Unix.waitpid [Unix.WUNTRACED] child_pid with
    | (_, Unix.WEXITED 0)           -> (VBool true,  env)
    | (_, Unix.WEXITED _)           -> (VBool false, env)
    | (dead_pid, Unix.WSIGNALED s)  ->
        Printf.eprintf "process %d interrupted by signal %d\n" dead_pid s;
        (VBool false, env)
    | (aslept_pid, Unix.WSTOPPED s) ->
        Printf.eprintf "process %d stopped by signal %d\n" aslept_pid s;
        (VBool false, env)
\end{lstlisting}

\begin{bloc}{Code retour et \texttt{VBool}}
Un code retour $0$ correspond à \texttt{VBool true} ;
tout autre code correspond à \texttt{VBool false}.
Ce comportement est \textbf{intentionnel} : des commandes comme \texttt{test}
retournent $1$ pour signifier « faux » --- ce n'est pas une erreur.
Le flag \texttt{Unix.WUNTRACED} (repris de \texttt{foo.ml}) permet de
détecter également les processus arrêtés (\texttt{Ctrl+Z}).
\end{bloc}

% ============================================================
\section{Pipes Unix}
% ============================================================

\subsection{Principe}

Inspirée de \texttt{pipe\_commands} dans \texttt{foo.ml},
\texttt{eval\_pipe} relie deux expressions par un tube Unix :

\[
\irule{
\sem{e_1}{\sigma}{\_}{\_}
\quad
\sem{e_2}{\sigma}{v}{\sigma'}
}{
\sem{\V{EPipe}(e_1,e_2)}{\sigma}{v}{\sigma'}
}
\]

\subsection{Protocole}

\begin{enumerate}
  \item \texttt{Unix.pipe()} crée la paire \texttt{(fd\_in, fd\_out)}.
  \item \texttt{stdin} est sauvegardé avec \texttt{Unix.dup(stdin)}.
  \item \textbf{Fils gauche} : redirige \texttt{stdout} vers \texttt{fd\_out},
        évalue $e_1$, puis quitte (\texttt{exit 0}).
  \item \textbf{Père} : redirige \texttt{stdin} depuis \texttt{fd\_in},
        évalue $e_2$.
  \item Le père attend la fin du fils (\texttt{waitpid}),
        puis restaure \texttt{stdin} original.
\end{enumerate}

\begin{lstlisting}[language=ml]
let (fd_in, fd_out) = Unix.pipe () in
let saved_stdin = Unix.dup Unix.stdin in
match Unix.fork () with
| 0 ->
    (* Fils gauche : ecrit dans le pipe *)
    Unix.dup2 fd_out Unix.stdout;
    Unix.close fd_out; Unix.close fd_in; Unix.close saved_stdin;
    ignore (eval e1 env); exit 0
| left_pid ->
    (* Pere : lit depuis le pipe *)
    Unix.dup2 fd_in Unix.stdin;
    Unix.close fd_out; Unix.close fd_in;
    let result = eval e2 env in
    ignore (Unix.waitpid [Unix.WUNTRACED] left_pid);
    Unix.dup2 saved_stdin Unix.stdin;
    Unix.close saved_stdin;
    result
\end{lstlisting}

\begin{bloc}{Différence avec \texttt{foo.ml}}
Dans \texttt{foo.ml}, le pipe est définitif : le processus père
n'a plus besoin de lire depuis le clavier après le pipe.
Dans CaioSHell, le père est le REPL lui-même.
Après l'exécution du pipe, il doit continuer à lire les commandes
de l'utilisateur.
\texttt{Unix.dup(stdin)} sauvegarde le descripteur original,
et \texttt{Unix.dup2(saved\_stdin, stdin)} le restaure
une fois le pipe terminé.
\end{bloc}

% ============================================================
\section{Résumé}
% ============================================================

\begin{center}
\begin{tabular}{|l|c|c|}
\hline
Construction & Étend $\sigma$ & Visible à l'appelant \\
\hline
\texttt{EDefVar}, \texttt{EAssign} & oui & oui \\
\texttt{EDefFun} & oui & oui \\
\texttt{EApp} (corps) & local & non \\
\texttt{ESeq} & propagation & oui \\
\texttt{EWhile}, \texttt{EFor} & propagation & oui \\
\texttt{EIf}, \texttt{EAnd} & branche exécutée & oui \\
\texttt{ECmd}, \texttt{EPipe} & non & non \\
littéraux, \texttt{EPass} & non & non \\
\hline
\end{tabular}
\end{center}

\begin{bloc}{Conclusion}
Le modèle sémantique :

\[
\sem{e}{\sigma}{v}{\sigma'}
\]

reflète directement la signature OCaml :

\begin{lstlisting}[language=ml]
eval : expr -> env -> value * env
\end{lstlisting}

Le projet combine :

\begin{itemize}
  \item interprétation dynamique sans phase de typage,
  \item portée lexicale via les fermetures (\texttt{VClosure}),
  \item appels système Unix (\texttt{fork}, \texttt{execvp}, \texttt{pipe}),
  \item propagation fonctionnelle de l'environnement,
  \item fonctions utilitaires réimplémentées sans le module \texttt{List}.
\end{itemize}
\end{bloc}

\begin{bloc}{Exemple de code testé}

\begin{lstlisting}
#>> CaioSHell - Tests

#>> Simple commands
echo " 1. Simple commands ";;
echo "Hello, CaioSHell!";;
pwd;;

#>> Variables
echo " Variables ";;
define var name = "Alice";;
echo @name;;
define var n = 42;;
echo @n;;

#>> Arithmetic
echo " Arithmetic ";;
define var a = 10;;
define var b = 3;;
define var sum  = @a + @b;;
define var diff = @a - @b;;
define var prod = @a * @b;;
define var quot = @a # @b;;
echo @sum;;
echo @diff;;
echo @prod;;
echo @quot;;

#>> String concatenation
echo " String concatenation ";;
define var s1 = "Hello, ";;
define var s2 = "World!";;
define var msg = @s1 + @s2;;
echo @msg;;

#>> Function definition and call
echo " Functions ";;
define function greet(person) {
    echo @person;
};;
greet("Alice");;
greet("Bob");;

#>> Function with arithmetic
echo " Function with arithmetic ";;
define function double(x) {
    define var result = @x * 2;
    echo @result;
};;
double(5);;
double(21);;

#>> Pipes
echo " Pipes ";;
echo "hello world" | wc -w;;
ls / | grep tmp;;

#>> And (&&)
echo " And (&&) ";;
(echo "first") && (echo "second");;

#>> If / else
echo "If / else ";;
if (test 1 "-eq" 1) {
    echo "condition true";
} else {
    echo "condition false";
};;
if (test 1 "-eq" 2) {
    echo "should not print";
} else {
    echo "correctly false";
};;

#>> For loop
echo " For loop ";;
for (@item in ["alpha", "beta", "gamma"]) {
    echo @item;
};;

#>> While loop (avec entiers negatifs)
echo " While loop ";;
define var i = 0;;
while (test @i "-gt" -3) {
    echo @i;
    i = @i - 1;
};;

#>> Variable reassignment
echo " Variable reassignment ";;
define var counter = 100;;
counter = @counter + 1;;
echo @counter;;

#>> Pass
echo " Pass ";;
if (test 0 "-eq" 1) {
    echo "unreachable";
} else {
    pass;
};;
echo "pass ok";;

#>> Files with dots in name
echo "File paths with dots ";;
cat /etc/hostname;;

#>> Flags (dash args)
echo " Command flags ";;
ls -l /tmp;;

#>> cd
echo "cd ";;
cd /tmp;;
pwd;;
cd /;;
pwd;;

echo " All tests done ";;
\end{lstlisting}

\end{bloc}

\end{document}
